% !TeX program = xelatex
\documentclass[UTF8,a4paper,11pt]{ctexart}

\usepackage[margin=2.0cm]{geometry}
\usepackage{amsmath,amssymb}
\usepackage{booktabs}
\usepackage{longtable}
\usepackage{array}
\usepackage{enumitem}
\usepackage{xcolor}
\usepackage{hyperref}
\usepackage{listings}
\usepackage{ragged2e}
\usepackage{microtype}

\setCJKmainfont{Microsoft YaHei}
\setCJKsansfont{Microsoft YaHei}
\setmonofont{Consolas}
\setCJKmonofont{Microsoft YaHei}

\hypersetup{colorlinks=true,linkcolor=blue!60!black,urlcolor=blue!60!black}
\lstset{basicstyle=\ttfamily\small,breaklines=true,frame=single,columns=fullflexible}
\setlist[itemize]{leftmargin=2em}
\setlist[enumerate]{leftmargin=2em}
\renewcommand{\arraystretch}{1.18}
\setlength{\LTleft}{0pt}
\setlength{\LTright}{0pt}
\setlength{\parskip}{0.35em}
\emergencystretch=3em
\sloppy
\newcolumntype{L}[1]{>{\RaggedRight\arraybackslash}p{#1}}
\newcommand{\code}[1]{\nolinkurl{#1}}

\title{中小计划排级列生成算法说明\\
\large \texttt{medium\_small\_raw} 实现}
\author{}
\date{\today}

\begin{document}
\maketitle
\tableofcontents
\newpage

\section{业务定位}

\code{medium_small_raw} 是一套中小计划排级列生成程序。它在大计划给出的航次--流向--箱区--尺寸配额基础上，生成两类结果：
\begin{itemize}
  \item \code{medium_plan.csv}：中计划输出，按粗属性组在箱区/贝位上的计划量汇总。
  \item \code{small_plan.csv}：小计划输出，按细属性组落到具体箱区、贝位和排号。
\end{itemize}

这套程序的核心目标有三项：
\begin{itemize}
  \item 细属性组尽量集中，减少同一细属性组跨箱区、跨 6 小贝区块和跨贝位。
  \item 粗属性组尽量箱区集中。小数量粗属性组倾向集中在一个箱区；大数量粗属性组允许使用多个箱区，但避免很小的尾巴箱区，并保持主要箱区之间相对均衡。
  \item 尽量继承大计划。箱量优先落在大计划给定的同航次、同流向、同尺寸箱区内；必要时按阶段逐步放宽箱区范围。
\end{itemize}

排级约束是本目录与贝位级列生成版本的主要区别。每个被选择的列都带有 \code{row_allocation}，小计划输出包含 \code{row_no}。模型不仅控制贝位总容量，还控制每个贝位每一排的物理容量、尺寸容量和不混排属性。默认排级不混属性为卸港 \code{port}，默认贝位级不混属性为尺寸 \code{size} 和箱高 \code{height}；这些属性可以通过输入中的 \code{attribute_rules} 覆盖。

\section{运行入口与输出}

\subsection{联合生成中计划和小计划}

联合生成入口为：
\begin{lstlisting}
.\.venv\Scripts\python.exe medium_small_raw\run_column_generation_planner.py
\end{lstlisting}

默认数据集为 \code{0519}。数据目录、大计划文件、目标航次、规划时刻和规划窗口都可以通过命令行参数覆盖。输出目录包含：
\begin{itemize}
  \item \code{small_plan.csv}：排级小计划。
  \item \code{medium_plan.csv}：由被选列汇总得到的中计划。
  \item \code{small_plan_six_bay_blocks.csv}：按连续 6 小贝区块汇总的小计划。
  \item \code{big_plan_used.csv}：本次使用的大计划配额行。
  \item \code{generated_columns.csv}：本次生成过的候选列。
  \item \code{diagnostics.json}：需求、列池、阶段、未放置箱、集中性、继承性和运行时间诊断。
\end{itemize}

\subsection{给定外部中计划后生成小计划}
\label{sec:external-medium}

外部中计划入口为：
\begin{lstlisting}
.\.venv\Scripts\python.exe medium_small_raw\run_small_plan_from_medium.py --medium-plan path\to\medium_plan.csv
\end{lstlisting}

该入口固定使用 \code{doc-only} 口径：只规划当前未进场资料箱，不生成预测兜底组。外部中计划形成两类硬上界：
\begin{itemize}
  \item 粗属性组--箱区上界 \((v,f,p,s,a)\)。
  \item 若外部中计划带有贝位字段，则形成粗属性组--箱区--贝位上界 \((v,f,p,s,a,b)\)。
\end{itemize}

小计划在这些上界、排级容量和不混排规则内落箱。若外部中计划给出的箱区/贝位本身无法容纳某些资料箱，未放置箱量会写入 \code{diagnostics.json}。

\section{输入口径}

\subsection{大计划}

大计划提供航次、流向、箱区、尺寸和计划箱量。当前读取器支持常见字段别名，例如 \code{voy_id}、\code{voyage_id}、\code{area_no}、\code{size}、\code{planned_qty}、\code{planned_boxes} 等。45ft 在继承大计划配额时归入 40ft 大计划尺寸口径；在小计划落排时仍按 45ft 规则处理。

大计划用于三件事：
\begin{itemize}
  \item 推导默认目标航次。没有显式传入 \code{--voyages} 时，程序读取规划日期上的 \code{OF} 行作为目标航次来源。
  \item 构造箱区配额上界 \(Q_{vfa\bar{s}}\)，其中 \(\bar{s}\) 为大计划尺寸口径。
  \item 形成继承偏好。精确大计划箱区成本最低，偏离大计划箱区会进入 fallback 层级成本和大计划模式偏离目标。
\end{itemize}

\subsection{堆场快照}

堆场快照提供箱区、贝位、排号和在场箱信息。程序从中构造：
\begin{itemize}
  \item 贝位物理容量和尺寸容量。
  \item 每个贝位每一排的物理容量和尺寸容量。
  \item 40ft/45ft 占用的连续小贝 footprint。
  \item 已有箱尺寸、箱高、卸港等属性。
  \item 当前已经在场的箱，用于从资料箱需求中剔除。
\end{itemize}

当某个贝位已有箱属性与待放需求组冲突时，该贝位或对应排不进入可行候选。排级可行性直接在列生成、整数主问题和 repair 中共同生效。

\subsection{箱区功能、船期和泊位距离}

箱区功能表用于判断某箱区是否支持某流向。船期和规划时刻用于计算目标航次窗口、需求比例和作业窗口。泊位距离表用于给航次--箱区组合增加距离成本。

箱区内其它装船作业也会进入目标：规划窗口内存在作业冲突会增加成本；规划窗口后 24 小时内有后续装船作业时，相关箱区会获得一定偏好。

\section{需求组}

\subsection{粗属性组}

粗属性组是中计划的业务单位，默认键为：
\[
g=(v,f,p,s),
\]
其中 \(v\) 为航次，\(f\) 为流向，\(p\) 为卸港，\(s\) 为真实尺寸 20、40 或 45。粗属性组箱量来自 \code{calculate_medium_demands}，并受大计划总配额约束。

\subsection{细属性组}

细属性组是小计划落排单位，默认键为：
\[
h=(v,f,p,s,\eta,w,\sigma,r),
\]
其中 \(\eta\) 为箱高，\(w\) 为重量等级，\(\sigma\) 为特殊积载代码，\(r\) 为预配/预留标记。资料箱读取后会先剔除当前已经在场的箱，再按细属性聚合。

\subsection{预测兜底组}

联合生成默认 \code{original} 模式下，模型内部还包含预测兜底组。预测兜底组用于补足中计划需求中尚未被资料箱覆盖的箱量，并参与容量占用和中计划输出。默认 \code{small_plan.csv} 只写出资料箱细属性组，预测兜底组不写入小计划。

需求模式含义如下：
\begin{itemize}
  \item \code{original}：中计划按中计划需求输出，小计划只输出资料箱；预测兜底组参与模型。
  \item \code{medium}：小计划输出包含资料箱和预测兜底组，总量对齐中计划需求。
  \item \code{medium-with-doc-floor}：保留全部资料箱，并补足中计划不足部分。
  \item \code{doc-only}：只规划当前未进场资料箱，不生成预测兜底组。
\end{itemize}

\section{排级列}

一列表示一个细属性组在一个贝位上的一批箱量放置方案：
\[
c=(h,b,q,R_c),
\]
其中 \(h\) 为细属性组，\(b\) 为贝位，\(q\) 为箱量，\(R_c\) 为排级分配。对 20ft 箱，footprint 是一个小贝；对 40ft/45ft 箱，footprint 是一对连续小贝。40ft/45ft 的 \code{row_allocation} 会同时记录 footprint 内两个小贝在同一排上的占用。

列生成时，候选列必须满足：
\begin{itemize}
  \item 箱区支持对应流向，或由用户箱区规则强制允许。
  \item 20ft 不放箱区边贝，45ft 只放箱区边贝。
  \item 40ft/45ft 必须存在可用的大贝 partner。
  \item 贝位已有尺寸、箱高等贝位级属性与需求组相容。
  \item 排已有卸港等排级属性与需求组相容。
  \item footprint 内所有小贝在同一排上都有足够的物理容量和尺寸容量。
  \item 列箱量不超过需求组剩余箱量、贝位容量、排容量和相关箱区配额。
\end{itemize}

每个候选贝位会生成若干箱量选项，既包含最大可放量，也包含小箱量和余数箱量，用于在集中和可行之间取得平衡。每个列在 \code{generated_columns.csv} 中写出 \code{row_allocation} 和 \code{stack_units}。

\section{四阶段业务流程}

求解过程采用“stage0 主分配 + stage1--stage3 repair 补齐 + repair LNS + 箱区内贝位重排”的流程。stage0 和 repair 决定箱区层方案；箱区内贝位重排在箱区层方案固定后运行，只重新安排箱区内部的贝位和排级可行列。

\begin{longtable}{L{0.13\linewidth}L{0.22\linewidth}L{0.55\linewidth}}
\toprule
阶段 & 业务含义 & 箱区范围和约束\\
\midrule
\endhead
stage0 & 大计划范围内的主体分配 & 候选箱区必须是同航次、同流向、同大计划尺寸口径下有精确大计划配额的箱区。大计划配额作为硬上界。列生成和主问题主要在这一范围内寻找继承性好、集中性好的方案。\\
stage1a & 精确大计划配额内补箱 & repair 中继续使用同 stage0 的精确大计划箱区范围，并严格执行大计划配额上界。\\
stage1b & 同航次流向的大计划箱区补箱 & 候选箱区放宽到该需求组所属航次、流向的大计划覆盖箱区，不再用精确箱区配额卡住单个需求组。\\
stage2 & 大计划覆盖箱区补箱 & 候选箱区放宽到任意被大计划覆盖的箱区，用于处理同航次流向范围内仍无法放下的箱。\\
stage3 & 功能兼容兜底补箱 & 候选箱区放宽到功能兼容且满足用户箱区策略、容量和排级规则的箱区。该阶段优先保证小计划可执行。\\
\bottomrule
\end{longtable}

每个阶段都遵守排级容量、不混排、边贝、尺寸、箱高、外部中计划配额等硬约束。阶段放宽的是箱区候选范围和大计划精确配额，不放宽物理和排级可行性。

stage0 后，程序先求一个“最少未放置箱”的整数模型。若该模型达到最优，则得到 \code{stage0_unplaced_cap}，后续最终主问题和 repair 都不能超过这个未放置箱上限。若 stage0 已经证明零未放置，后续阶段保持零未放置。

\section{主问题和硬约束}

主问题选择当前列池中的列。整数阶段每列为二进制变量 \(x_c\)，每个细属性组有未放置变量 \(u_h\)。需求覆盖约束为：
\[
\sum_{c:h(c)=h} q_cx_c+u_h=d_h,\quad \forall h.
\]

贝位、排和配额硬约束包括：
\[
\sum_{c:b(c)=b} q_cx_c \le Cap_b,
\]
\[
\sum_{c:b(c)=b,s(c)=s} q_cx_c \le Cap_{bs},
\]
\[
\sum_{c:(b,r)\in R_c} q_{cr}x_c \le Cap_{br},
\]
\[
\sum_{c:(b,r,s)\in R_c} q_{cr}x_c \le Cap_{brs},
\]
\[
\sum_{c:(v,f,a,\bar{s})} q_cx_c \le Q_{vfa\bar{s}}.
\]

对于同一贝位，模型用选择变量限制其尺寸、箱高等贝位级属性只能取一个值；对于同一排，模型限制卸港等排级属性只能取一个值。40ft/45ft 的 footprint 约束保证两边小贝同步占用，排级分配也保证两边小贝在对应排都有容量。

外部中计划入口还会加入：
\[
\sum_{c:(v,f,p,s,a)}q_cx_c\le E_{vfpsa},
\]
以及可选的贝位级中计划上界：
\[
\sum_{c:(v,f,p,s,a,b)}q_cx_c\le E_{vfpsab}.
\]

未放置变量用于表达当前硬约束下无法覆盖的剩余箱量。它不是容量或配额的替代品；容量、配额和排级规则始终是硬约束。最终方案优先比较未放置箱数，其次比较目标能量，最后比较选中列数。

\section{列生成、diving 和 repair}

列生成从初始列池开始，反复求解线性松弛主问题，并根据对偶价格对未进入列池的候选列定价。负检验数列优先进入列池。若常规定价停滞，程序还会加入低检验数列和整数可行性导向的 primal expansion 列，使后续整数主问题拥有更丰富的选择。

列生成结束后，程序在当前列池上执行 diving 主问题。diving 逐步固定 LP 解中的高价值列，并在遇到 LP 未放置箱时收缩固定批次或停止。diving 后还会执行若干局部 improvement 轮，释放部分细属性组做小规模重优化。

如果主问题结果仍有未放置箱，程序进入未放置修复：
\begin{itemize}
  \item staged greedy repair 按 stage1a、stage1b、stage2、stage3 顺序补箱，并比较默认组顺序和粗属性组集中优先顺序。
  \item repair LNS 以当前最好解为 incumbent，每轮释放一批细属性组，在剩余容量和剩余配额内求小规模 SCIP 子问题。
  \item 粗属性组 compaction LNS 默认关闭，打开后用于额外压实粗属性组箱区分布。
\end{itemize}

repair 候选方案按如下顺序排序：
\[
(\text{未放置箱数},\ \text{目标能量},\ \text{选中列数}).
\]
因此 repair LNS 只有在未放置箱数更少，或未放置箱数相同但目标更好时才替换 incumbent。零未放置 incumbent 下，LNS 子问题保持零未放置，不通过多留箱来换集中性。

repair 结束后，程序执行箱区内贝位重排。该步骤固定每个细属性组在每个箱区的箱量 \((group\_id, area\_no)\)，不改变箱区层分配、不改变大计划继承比例，也不改变每个粗属性组在每个箱区的合计箱量。重排只在同一箱区内重新选择贝位，并为所选贝位生成排级可行的 \code{row_allocation}。排号在这里主要用于证明贝位方案能够真实落排，重排评分只关心细属性组贝位集中、粗属性组箱区内贝位集中和小尾巴贝位减少。

箱区内贝位重排逐箱区尝试。某个箱区若无法在排级硬约束下完成重排，则该箱区保留 repair 后的原贝位方案；其它箱区仍可继续尝试。只有当整体贝位集中性评分变好，且所有 \((group\_id, area\_no)\) 箱量完全一致时，重排方案才会成为最终方案。若外部中计划给定贝位级硬配额，箱区内贝位重排不运行，因为贝位层方案已经由外部输入固定。

\section{目标函数}

\subsection{细属性组集中}

同一细属性组跨箱区、跨 6 小贝区块、跨贝位都会产生惩罚。对应参数为：
\begin{itemize}
  \item \code{--fine-group-area-penalty}，默认 80。
  \item \code{--fine-group-block-penalty}，默认 35。
  \item \code{--fine-group-bay-penalty}，默认 8。
\end{itemize}

\subsection{粗属性组集中}

同一粗属性组在同一箱区内使用更多 6 小贝区块和更多贝位会产生惩罚，参数为 \code{--coarse-area-block-penalty} 和 \code{--coarse-area-bay-penalty}。

粗属性组箱区目标分两类：
\begin{itemize}
  \item 小数量粗属性组：需求量不超过 \code{--medium-concentrated-group-threshold}，默认 26。目标强烈惩罚跨箱区，并奖励最大箱区承载更多箱量。
  \item 大数量粗属性组：不强求只用极少箱区，而是惩罚低于 \code{--medium-large-group-min-area-boxes} 的小尾巴箱区，默认阈值 10；同时用均衡项避免主要箱区之间差异过大。
\end{itemize}

大数量粗属性组的相关参数包括 \code{--medium-large-group-small-area-penalty}、\code{--medium-large-group-area-open-penalty} 和 \code{--medium-large-group-target-area-boxes}。

箱区内贝位重排阶段继续使用“细属性组少跨贝位、粗属性组少跨贝位、减少小尾巴贝位”的集中性口径。该阶段不再考虑大计划偏离、泊位距离、作业窗口等目标，因为箱区层方案已经固定；它只在不改变箱区量且满足排级可行性的前提下改善贝位层分布。

\subsection{大计划继承}

大计划继承由三部分表达：
\begin{itemize}
  \item 精确大计划箱区配额上界。stage0 和 stage1a 严格执行同航次、同流向、同大计划尺寸口径的箱区配额。
  \item \code{--required-area-reward}。列落在大计划要求箱区时获得基础成本奖励，默认 1000。
  \item \code{--big-plan-fallback-tier-penalty} 和 \code{--big-plan-area-deviation-penalty}。前者按候选箱区偏离大计划的层级增加成本；后者惩罚航次--流向--大计划尺寸合计分布相对大计划比例的偏离。
\end{itemize}

大计划偏离目标约束的是合计箱区模式，不要求每个粗属性组都机械地按大计划比例拆分。粗属性组集中目标和大计划继承目标共同决定最终箱区选择。

\subsection{泊位距离和作业窗口}

航次--箱区使用会带来泊位距离成本。规划窗口内已有其它装船作业的箱区成本更高；规划窗口后 24 小时内有后续装船作业的箱区会有一定偏好。stage3 选择大计划之外的功能兜底箱区时，这类成本有助于把箱放在更合理的作业位置。

\section{排级可行性的含义}

排级可行性表示：对 \code{small_plan.csv} 中每一行，输出的 \code{area_no}、\code{bay_no}、\code{row_no} 和 \code{planned_boxes} 能映射回列的 \code{row_allocation}；同一贝位同一排上的箱量不超过该排物理容量和尺寸容量；同一排不会出现不允许混放的属性组合。

对 40ft/45ft 箱，\code{row_allocation} 记录 footprint 内两个小贝的排级占用。只有当两个小贝在同一排上都满足容量和属性相容时，该列才是可行列。这样可以避免“贝位层面看似可行，落到排后无法摆放”的情况。

\section{中计划和小计划的一致性}

联合生成时，\code{medium_plan.csv} 直接由被选列按粗属性组、箱区和贝位汇总得到。每一行包含：
\begin{flushleft}\small\ttfamily
plan\_level, voyage\_id, flow, port, size, area\_no, bay\_no, six\_bay\_block\_id, six\_bay\_block\_bays, planned\_boxes, document\_boxes, forecast\_fallback\_boxes
\end{flushleft}

\code{small_plan.csv} 则进一步落到排，包含：
\begin{flushleft}\small\ttfamily
plan\_level, voyage\_id, group\_id, demand\_source, flow, port, size, height, weight\_class, special\_stow, special\_stow\_code, area\_no, bay\_no, row\_no, row\_allocation, six\_bay\_block\_id, six\_bay\_block\_bays, six\_bay\_block\_total\_boxes, planned\_boxes
\end{flushleft}

默认 \code{original} 口径下，小计划只写资料箱，预测兜底组只体现在中计划的 \code{forecast_fallback_boxes} 中。诊断文件中的 \code{small_medium_consistency_violations} 和 \code{small_medium_bay_consistency_violations} 用于检查小计划汇总是否超过中计划对应口径。

\section{诊断字段}

\code{diagnostics.json} 是判断结果是否正常的主要文件。常用字段包括：
\begin{itemize}
  \item \code{planned_medium_boxes}、\code{planned_small_boxes}、\code{unplaced_boxes}。
  \item \code{stage0_unplaced_cap} 和 \code{stage0_unplaced_cap_source}。
  \item \code{master_status}、\code{master_algorithm}、\code{final_column_count}、\code{selected_column_count}。
  \item \code{repair_lns_rounds_run}、\code{repair_lns_improvements}、\code{repair_lns_stop_reason}。
  \item \code{post_repair_area_relayout_enabled}、\code{post_repair_area_relayout_accepted}、\code{post_repair_area_relayout_before}、\code{post_repair_area_relayout_after}，记录箱区内贝位重排是否运行、是否接受以及重排前后的贝位集中性指标。
  \item \code{medium_big_plan_inheritance.inheritance_ratio}。
  \item \code{medium_fragmentation} 下的 \code{tiny_area_rows}、\code{small_coarse_multi_area_groups}、\code{large_coarse_tiny_area_rows}。
  \item \code{attribute_rules}，记录本次使用的粗属性、细属性、贝位不混属性和排不混属性。
\end{itemize}

这些字段共同解释三类问题：是否所有箱都放下、粗细属性组集中性是否合理、大计划继承性是否保持在可接受范围内。

\section{主要参数}

\begin{longtable}{L{0.94\linewidth}}
\toprule
参数、默认值与业务含义\\
\midrule
\code{--max-iterations=30}：列生成最大迭代次数。\\
\code{--columns-per-iteration=2500}：每轮最多加入的负检验数列数量。\\
\code{--initial-columns-per-group=16}：每个细属性组初始候选贝位数量。\\
\code{--max-candidate-bays-per-group=500}：每个细属性组最多保留的候选贝位数量。\\
\code{--mip-time-limit=120}、\code{--mip-gap=0.01}：最终主问题求解限制。\\
\code{--diving-improvement-rounds=6}、\code{--diving-improvement-time-limit=6.0}：diving 后局部重优化轮数和单轮时间。\\
\code{--repair-lns-rounds=2}、\code{--repair-lns-time-limit=3.0}：未放置修复 LNS 轮数和单轮时间。\\
\code{--coarse-compaction-lns-rounds=0}：额外粗属性组压实 LNS，默认关闭。\\
\code{post\_repair\_area\_relayout\_enabled=True}：repair 后箱区内贝位重排，固定每个细属性组在每个箱区的箱量，只优化箱区内贝位集中性。\\
\code{post\_repair\_area\_relayout\_max\_patterns=1}：箱区内贝位重排时每个候选贝位只取一个排级可行模式，排只作为可行性证明，不作为主要优化目标。\\
\code{--medium-concentrated-group-threshold=26}：小数量粗属性组和大数量粗属性组的分界。\\
\code{--medium-small-group-area-split-penalty=2400}：小数量粗属性组跨箱区惩罚。\\
\code{--medium-small-group-fragment-penalty=90}：小数量粗属性组碎片惩罚。\\
\code{--medium-large-group-min-area-boxes=10}：大数量粗属性组小尾巴箱区阈值。\\
\code{--medium-large-group-small-area-penalty=900}：大数量粗属性组小尾巴箱区惩罚。\\
\code{--medium-large-group-target-area-boxes=60}：大数量粗属性组目标箱区数量的参考箱量。\\
\code{--unplaced-penalty=100000}：未放置变量目标权重；未放置控制主要由 stage0 上限和 repair 不恶化约束保证。\\
\code{--required-area-reward=1000}：落在大计划要求箱区的奖励。\\
\code{--big-plan-area-deviation-penalty=3}：大计划合计箱区模式偏离惩罚。\\
\code{--big-plan-fallback-tier-penalty=20}：偏离大计划候选层级的惩罚。\\
\code{--full-column-pool}：直接枚举当前可枚举列池，跳过常规定价循环。\\
\code{--no-scip}：使用确定性贪心回退，适合检查数据链路，不代表正式优化质量。\\
\bottomrule
\end{longtable}

\end{document}
